\documentclass[12pt,a4paper,twoside,openany]{book}

% Pacchetti essenziali
\usepackage[utf8]{inputenc}
\usepackage[T1]{fontenc}
\usepackage[english]{babel}
\usepackage{amsmath,amsfonts,amssymb}
\usepackage{graphicx}
\usepackage{float}
\usepackage[section]{placeins}
\usepackage{booktabs}
\usepackage{longtable}
\usepackage{array}
\usepackage{multirow}
\usepackage{multicol}
\usepackage{url}
\usepackage{hyperref}
\usepackage{cite}
\usepackage{algorithm}
\usepackage{algorithmic}
\usepackage{listings}
\usepackage{xcolor}
\usepackage{tikz}
\usetikzlibrary{shapes, arrows.meta, positioning, calc, chains, decorations.pathreplacing, shadows, fit, backgrounds}
% ==== LLM DIAGRAM LIBRARY ====
\usepackage{tikz}
\usetikzlibrary{positioning,arrows.meta,fit,shadows.blur,shapes.geometric,decorations.pathreplacing}

% Stili generali
\tikzset{
  llm/base/.style={
    draw,
    rounded corners=3pt,
    align=center,
    font=\sffamily\small,
    blur shadow={shadow blur steps=5}
  },
  llm/module/.style={
    llm/base,
    minimum width=2.8cm,
    minimum height=1.2cm,
    fill=blue!15,
    draw=blue!50!black,
    thick
  },
  llm/io/.style={
    llm/base,
    minimum width=2.5cm,
    minimum height=1cm,
    fill=gray!10,
    draw=gray!50!black,
    dashed
  },
  llm/arrow/.style={
    -Latex,
    thick,
    draw=black!80
  },
  llm/group/.style={
    draw,
    dashed,
    inner sep=8pt,
    rounded corners=5pt,
    fill=blue!5,
    draw=blue!30
  }
}

% Ambiente comodo per i diagrammi di architettura
\newenvironment{llmarchitecture}
{%
  \begin{tikzpicture}[
    node distance=1.5cm and 1.5cm,
    every node/.style={font=\sffamily}
  ]%
}
{%
  \end{tikzpicture}%
}

\usepackage{pgfplots}
\usepackage{subcaption}
\usepackage{geometry}
\usepackage{fancyhdr}
\usepackage{setspace}

% Configurazione della pagina
\geometry{
    portrait,
    left=3cm,
    right=2.5cm,
    top=3cm,
    bottom=3cm,
    bindingoffset=0.5cm
}

% Configurazione dell'interlinea
\onehalfspacing

% Configurazione dei link
\hypersetup{
    colorlinks=true,
    linkcolor=black,
    filecolor=magenta,
    urlcolor=blue,
    citecolor=blue,
    pdftitle={SVG Captioning with Transformer Models: Advanced Fine-tuning Techniques - PART 1},
    pdfauthor={Emanuele Di Luzio},
    pdfsubject={Master's Thesis - Chapters 1-3},
}

% Configurazione del codice
\lstset{
    basicstyle=\ttfamily\footnotesize,
    breaklines=true,
    frame=single,
    numbers=left,
    numberstyle=\tiny,
    showstringspaces=false,
    commentstyle=\color{gray},
    keywordstyle=\color{blue},
    stringstyle=\color{red}
}

% Configurazione delle intestazioni
\pagestyle{fancy}
\fancyhf{}
\fancyhead[LE,RO]{\thepage}
\fancyhead[LO]{\rightmark}
\fancyhead[RE]{\leftmark}
\renewcommand{\headrulewidth}{0.4pt}
\setlength{\headheight}{14.5pt}

\begin{document}

% Frontespizio
\begin{titlepage}
    \centering
    \vspace*{2cm}
    
    {\Large\textbf{UNIVERSITY OF MODENA AND REGGIO EMILIA}}\\[0.5cm]
    {\large Department of Engineering "Enzo Ferrari"}\\[0.5cm]
    {\large Master's Degree in Computer Engineering}\\[3cm]
    
    {\Huge\textbf{SVG Captioning with Transformer Models: Advanced Fine-tuning Techniques}}\\[0.5cm]
    {\Large\textit{PART 1: Chapters 1-3}}
    \vspace{2cm}
    \begin{minipage}{0.4\textwidth}
        \begin{flushleft}
            \large
            \textbf{Candidate:}\\
            Emanuele Di Luzio\\
            Student ID: \rule{4cm}{0.4pt}
        \end{flushleft}
    \end{minipage}
    \hfill
    \begin{minipage}{0.4\textwidth}
        \begin{flushright}
            \large
            \textbf{Supervisor:}\\
            Prof.\ \rule{5cm}{0.4pt}\\
            \vspace{0.5cm}
            \textbf{Co-supervisor:}\\
            Dr.\ \rule{5cm}{0.4pt}
        \end{flushright}
    \end{minipage}
    
    \vfill
    {\large Academic Year 2023-2024}
\end{titlepage}

% Indice
\tableofcontents
\newpage

% Abstract
\chapter*{Abstract}
\addcontentsline{toc}{chapter}{Abstract}

Scalable Vector Graphics (SVG) are ubiquitous in modern web design, yet current Multimodal Large Language Models (MLLMs) struggle to interpret their raw XML structure. Most existing approaches rely on rasterizing SVGs into pixel grids, a process that discards the semantic richness of the vector definition and introduces resolution artifacts.

In this work, I present a novel architecture that enables a decoder-only LLM to directly ingest and interpret SVG code without rasterization. I integrate a \textbf{Spatial Position Encoder (SPE)}--originally developed by Zini \textit{et al.}--that maps continuous geometric coordinates into the LLM's embedding space using sinusoidal functions. By combining this encoder with \textbf{Qwen2-7B} and \textbf{Gemma-9B} via \textbf{Low-Rank Adaptation (LoRA)}, I demonstrate that the model can learn to "read" vector geometries effectively.

My experiments, conducted on a dataset of \textbf{20,000 SVG-caption pairs} and evaluated on a stratified benchmark of \textbf{400 samples}, show that this vector-native approach outperforms zero-shot raster baselines. Specifically, the \textbf{SPE+Qwen2-7B} configuration achieves a CLIPScore of 29.3 and a BLEU-1 score of 0.42, offering a parameter-efficient alternative to vision-encoder-based methods. I provide a detailed analysis of the model's capabilities and limitations, offering a new perspective on vector-language understanding.

\chapter*{Sommario}
\addcontentsline{toc}{chapter}{Sommario}

Le immagini vettoriali (SVG) sono onnipresenti nel design digitale, tuttavia gli attuali modelli multimodali (MLLM) faticano a interpretare la loro struttura XML grezza. La maggior parte degli approcci esistenti si basa sulla rasterizzazione degli SVG in griglie di pixel, un processo che scarta la ricchezza semantica della definizione vettoriale e introduce artefatti di risoluzione.

In questa tesi, presento un'architettura che consente a un LLM decoder-only di ingerire e interpretare direttamente il codice SVG senza rasterizzazione. Integro uno \textbf{Spatial Position Encoder (SPE)}--originariamente sviluppato da Zini \textit{et al.}--che mappa le coordinate geometriche continue nello spazio di embedding dell'LLM utilizzando funzioni sinusoidali. Combinando questo encoder con \textbf{Qwen2-7B} e \textbf{Gemma-9B} tramite \textbf{Low-Rank Adaptation (LoRA)}, dimostro che il modello può imparare a "leggere" le geometrie vettoriali in modo efficace.

I miei esperimenti, condotti su un dataset di \textbf{20.000 coppie SVG-caption} e valutati su un benchmark stratificato di \textbf{400 campioni}, mostrano che questo approccio vettoriale supera le baseline raster zero-shot. Nello specifico, la configurazione \textbf{SPE+Qwen2-7B} ottiene un CLIPScore di 29.3 e un punteggio BLEU-1 di 0.42, offrendo un'alternativa efficiente ai metodi basati su vision-encoder. Fornisco un'analisi dettagliata delle capacità e dei limiti del modello, offrendo una nuova prospettiva sulla comprensione del linguaggio vettoriale.

% ============================================
% CAPITOLI 1-3
% ============================================

\chapter{Introduction}
\label{chap:introduzione}

The interpretation of visual information by artificial intelligence has undergone a remarkable transformation over the past decade. From early Convolutional Neural Networks (CNNs) designed to classify static images, to modern Vision Transformers (ViTs) capable of understanding complex scenes, the field has consistently pushed the boundaries of what machines can ``see'' and comprehend. Yet, despite these advances, a fundamental assumption has remained largely unchallenged: that visual data consists of pixel grids--discrete arrays of color intensity values.

This pixel-centric paradigm, while successful for photographs and natural imagery, overlooks a critical dimension of modern digital content. A substantial portion of the visual material I encounter daily--icons, logos, technical diagrams, user interface elements, typographic designs--is not created from pixels at all. Instead, these graphics are defined as \textbf{Scalable Vector Graphics (SVG)}, mathematical descriptions of shapes, curves, and paths that are inherently resolution-independent and semantically structured.

Consider a simple icon: a magnifying glass representing ``search'' functionality on a website. To a human designer, this icon is not a collection of colored dots, but a composition of geometric primitives--a circle for the lens, a line for the handle, perhaps a smaller circle to suggest a reflection. The icon is encoded as a sequence of XML instructions that describe these shapes precisely:
\begin{itemize}
    \item \texttt{<circle cx="45" cy="45" r="35"/>} defines the lens;
    \item \texttt{<line x1="70" y1="70" x2="100" y2="100"/>} defines the handle.
\end{itemize}

\begin{figure}[ht]
    \centering
    \begin{tikzpicture}[scale=0.05, yscale=-1, every node/.style={scale=1}]
        % The Magnifying Glass Icon (Optimized Geometry)
        \draw[line width=3pt] (45,45) circle (35);
        \draw[line width=3pt, line cap=round] (70,70) -- (100,100);
    \end{tikzpicture}
    \caption{Visual representation of the simple SVG magnifying glass defined by the XML code.}
    \label{fig:intro_magnifying_glass}
\end{figure}
\noindent This representation is compact, editable, and infinitely scalable; zooming in never introduces pixelation or blur. However, when state-of-the-art Multimodal Large Language Models (MLLMs) --- systems like GPT-4 with Vision (GPT-4V), LLaVA, or BLIP-2 --- are asked to describe such an icon, they cannot read the vector definition directly. Instead, the standard practice is to \textbf{rasterize} the SVG: render it to a fixed-resolution bitmap (\textit{e.g.}, 512$\times$512 pixels), then feed this raster image to a vision encoder. This approach has two significant drawbacks. First, it discards the rich semantic structure embedded in the SVG markup--the hierarchy of groups, the annotations in \texttt{id} and \texttt{class} attributes, the explicit geometric relationships between shapes. Second, it introduces aliasing and discretization artifacts that can obscure fine details, especially for icons with thin lines or small features.

This thesis challenges the prevailing pixel-based paradigm by posing a simple question: \textit{Can I teach a language model to "read" vector graphics natively, without converting them to raster images first?} The motivation for this question is both practical and philosophical. On the practical side, there is an urgent need for automatic captioning systems that can generate accurate textual descriptions of digital graphics. On the philosophical side, treating SVGs as a \textit{language} aligns naturally with the strengths of Large Language Models (LLMs).

\section{The Challenge of Vector Graphics Captioning}
\label{sec:challenge-vector-captioning}

To appreciate the difficulty of the task, I must first understand what makes SVG fundamentally different from raster images. An SVG file is an XML document conforming to the W3C Scalable Vector Graphics specification \cite{w3c2011svg}. Unlike a JPEG or PNG, which stores a fixed grid of pixel values, an SVG stores a \textit{hierarchical DOM tree}--a structured set of graphical elements. Each element (\textit{e.g.}, \texttt{<path>}, \texttt{<circle>}, \texttt{<rect>}) is defined by a set of attributes: geometric parameters (coordinates, radii, control points), style properties (fill color, stroke width, opacity), and possibly metadata (titles, descriptions, ARIA labels).

The flexibility of this representation is a double-edged sword. On one hand, it enables powerful features like infinite zoom, easy editing, and style inheritance. On the other hand, it introduces significant variability. The same visual icon can be encoded in countless ways: paths can be absolute or relative, clockwise or counterclockwise; shapes can be grouped or flattened; transformations (translate, rotate, scale) can be applied at any level of the hierarchy. Furthermore, the coordinate values are \textit{continuous}--floating-point numbers that specify precise locations on a two-dimensional canvas. This continuous nature poses a challenge for standard text tokenization schemes, which typically operate on discrete vocabularies.

Moreover, the ``vocabulary'' of SVG includes geometric concepts that do not map straightforwardly to the semantic categories learned by vision-language models trained on natural images. A curve defined by a cubic Bézier path has no direct pixel-level analogue; it is a parametric function. Understanding that a specific sequence of path commands draws, say, a magnifying glass, requires integrating the geometric information across multiple elements and inferring the high-level semantic category from the emergent shape.

This gap between geometry and semantics is precisely where existing Multimodal Large Language Models (MLLMs) struggle. Standard vision encoders--whether ResNet-based or Vision Transformer-based--are designed to extract features from dense pixel grids. They excel at recognizing textures, edges, and spatial patterns in natural photos, but they have no mechanism to directly ingest the symbolic, structured representation of an SVG. Consequently, the common workaround is rasterization: render the SVG to a pixel image and process it as if it were a photograph. This approach, while expedient, is fundamentally mismatched to the data modality. It is analogous to asking someone to describe a written sentence by first printing it, photographing the page, and then performing OCR--technically feasible, but wasteful and error-prone.

\section{Accessibility, Search, and the Need for Automatic Captions}
\label{sec:motivation-accessibility}

The practical importance of SVG captioning extends far beyond academic curiosity. In the context of web accessibility, the World Wide Web Consortium (W3C) Web Content Accessibility Guidelines (WCAG) mandate that all non-text content be accompanied by alternative text (alt-text) that conveys the same information. For images, this is typically a short caption describing the visual content. However, many websites, design systems, and icon libraries contain thousands of SVG icons, and manually writing captions for each is a labor-intensive and often neglected task.

Consider a large design system like Google's Material Icons or Font Awesome, which provide hundreds of icons for user interface elements. Each icon needs a descriptive label not only for accessibility compliance but also for searchability. A designer looking for a ``notification bell'' icon should be able to search for that phrase and find the relevant SVG. Similarly, screen reader users navigating a web page should hear ``notification bell icon'' rather than the generic fallback ``image'' or, worse, the cryptic SVG filename ``icon\_287.svg''.

Current practice often relies on manually assigned tags or filenames, which are inconsistent and prone to error. An automatic captioning system that can generate accurate, human-readable descriptions directly from the vector definition would significantly streamline the design workflow, improve accessibility, and enable more sophisticated search and retrieval applications (\textit{e.g.}, semantic similarity search across icon libraries). Furthermore, as vector graphics are increasingly used in domains beyond traditional iconography--technical diagrams, data visualizations, UI mockups, schematic drawings--the ability to automatically annotate and describe these graphics becomes even more valuable. For instance, in educational contexts, automatically generating descriptions of geometric diagrams could aid visually impaired students. In software engineering, automatically documenting UI components from their SVG prototypes could improve code documentation and design handoff.

\section{Research Hypothesis: Treating SVG as a Structured Language}
\label{sec:hypothesis}

The central hypothesis of this thesis is that SVG captioning is best approached not as a vision problem (processing pixels), but as a \textbf{structured language understanding problem}. Since an SVG is fundamentally a symbolic representation--a sequence of XML tags and numeric coordinates--it lies closer to the domain of text than to the domain of natural images. This observation suggests that Large Language Models (LLMs), which are inherently designed to process and generate sequences of symbols, might be well-suited to this task.

However, a naive approach --- simply feeding the raw SVG XML code as a text string to an LLM --- fails in practice. The coordinate values are verbose (\textit{e.g.}, \texttt{d="M12.5,34.2..."}) and the numeric precision is lost during text tokenization. Early experiments with zero-shot prompting produced generic or hallucinatory captions.

To address this, I first investigate a \textbf{text-only fine-tuning baseline}, where the LLM is adapted to process linearized SVG code as a standard language modeling task. While this improves syntactic understanding, it remains limited in grasping complex geometries. Therefore, I propose a hybrid approach that bridges the gap between the continuous geometric domain and the discrete symbolic domain of language. Specifically, I integrate a \textbf{SVG Path Embedder (SPE)}--a neural module developed by Zini \textit{et al.} \cite{zini2025vhector}--that maps the continuous $(x, y)$ coordinates of SVG elements into a high-dimensional embedding space compatible with the LLM's input representation. This encoder does not operate in isolation; it must be integrated with a powerful language generation component that can transform the geometric features into fluent natural language, which is where Large Language Models enter the picture.

\section{The Role of Large Language Models in the Proposed Approach}
\label{sec:llms-role}

Why use Large Language Models at all for this task? The answer lies in their remarkable capacity for semantic understanding and text generation. LLMs like GPT-4, Llama, Qwen, and Gemma are pre-trained on vast corpora of text, giving them extensive world knowledge and the ability to generate fluent, contextually appropriate natural language. Crucially, these models are \textit{decoder-only} architectures, meaning they generate text autoregressively--one token at a time--conditioned on a prefix of input tokens.

In my architecture, the SPE produces a sequence of visual embeddings corresponding to the geometric elements of the SVG. These embeddings are then prepended (or interleaved) with the LLM's text token embeddings, and the model is tasked with generating a caption autoregressively. The LLM's pre-trained linguistic knowledge allows it to map the geometric features to semantic categories (\textit{e.g.}, recognizing that a circle with a line forms a ``magnifying glass''), and to generate grammatically correct, human-readable descriptions. However, off-the-shelf LLMs are not trained to interpret geometric embeddings. Their input space consists of text tokens (words, subwords, characters), not continuous vector representations of shapes. Therefore, I must \textit{adapt} the LLM to this new modality. This adaptation is achieved through fine-tuning: I train the model on a dataset of paired SVG-caption examples, teaching it to associate the geometric embeddings produced by the SPE with the corresponding textual descriptions.

\section{Parameter-Efficient Fine-Tuning with LoRA}
\label{sec:peft-intro}

Fine-tuning a modern LLM with billions of parameters is computationally expensive and risks \textit{catastrophic forgetting}--the model may overfit to the new task and lose its general linguistic capabilities. Full fine-tuning also requires storing a separate copy of the entire model for each task, which is impractical when deploying multiple specialized systems.

To mitigate these issues, I employ \textbf{Parameter-Efficient Fine-Tuning (PEFT)}, specifically \textbf{Low-Rank Adaptation (LoRA)} \cite{hu2021lora}. LoRA freezes the pre-trained weights of the LLM and introduces small, trainable low-rank matrices that modify the model's behavior. This drastically reduces the number of trainable parameters (from billions to millions) and accelerates training, while preserving the core linguistic knowledge of the base model. In my experiments, I apply LoRA to the attention projection matrices (query, key, value, and output projections) of the Transformer layers. During training, only the LoRA matrices and the SPE parameters are updated; the LLM's pre-trained weights remain frozen. This not only makes training feasible on limited hardware (\textit{e.g.}, a single or pair of GPUs), but also ensures that the adapted model retains its ability to generate fluent, grammatically correct captions.

\section{The SVG Path Embedder: Bridging Geometry and Language}
\label{sec:spe-intro}

With the LLM and the fine-tuning strategy established, I arrive at the final piece of the puzzle: the SVG Path Embedder (SPE). The SPE is a custom neural module designed to convert the continuous geometric properties of SVG elements into dense vector representations that the LLM can process. I provide a high-level preview here; the full technical details, including design alternatives and ablation studies, are developed in Chapter~\ref{chap:metodologia}.

The core challenge is encoding continuous coordinates. Unlike discrete token indices in text (which have a natural ordering and are easily represented as one-hot vectors or learned embeddings), $(x, y)$ coordinates are real-valued and do not fit naturally into the vocabulary of a language model. Naive discretization (\textit{e.g.}, rounding to the nearest integer or binning into fixed ranges) introduces quantization artifacts and limits the model's ability to perceive fine geometric details.

Instead, the SPE employs a technique inspired by the sinusoidal positional encodings used in the original Transformer architecture and by Fourier feature mappings in Neural Radiance Fields (NeRF). The core idea is to project each coordinate value through a bank of sinusoidal functions with geometrically increasing frequencies:
\[
\text{PE}(x) = \left[ \sin(\omega_1 x), \cos(\omega_1 x), \sin(\omega_2 x), \cos(\omega_2 x), \ldots, \sin(\omega_k x), \cos(\omega_k x) \right]^T
\]
where $\omega_k = 10000^{2k/d}$ for $k = 0, 1, \ldots, d/2 - 1$. This projection maps each scalar coordinate to a high-dimensional vector. The lower frequencies capture coarse, global position information (``is this shape in the top-left or bottom-right of the canvas?''), while the higher frequencies encode fine-grained details (``is this corner sharp or slightly rounded?''). By concatenating these positional encodings with learnable embeddings for element type (circle, path, rectangle, etc.) and style attributes (fill, stroke, etc.), and passing the result through a multi-layer perceptron (MLP), the SPE produces a unified embedding for each SVG element. These embeddings are then linearly projected to match the dimensionality of the LLM's input space and normalized to ensure numerical stability during training.

\section{Research Objectives}
\label{sec:research-objectives}

The primary objective of this thesis is to develop and validate a system capable of generating accurate, human-readable textual captions for SVG graphics by processing their vector definition directly, without rasterization. This broad goal is decomposed into four specific research objectives: (i) to \textbf{develop and validate a Baseline Training Process (LoRA-LLM)}, implementing a text-only fine-tuning pipeline where the SVG code is treated as a raw string to establish a strong baseline; (ii) to \textbf{develop and validate a Multimodal Training Process (SPE-LoRA)}, integrating the \textbf{SVG Path Embedder (SPE)} with decoder-only LLMs via parameter-efficient fine-tuning to enable true geometric understanding; (iii) to \textbf{curate a rigorous evaluation benchmark} consisting of a diverse, stratified test set of SVG-caption pairs; and (iv) to \textbf{conduct quantitative and qualitative analysis} of the model's performance, identifying strengths and failure modes.



\section{Summary of Methodology and Key Findings}
\label{sec:summary-methodology}

This thesis introduces a vector-native SVG captioning pipeline centered on the SVG Path Embedder (SPE), which maps geometric and structural properties of SVG elements into a continuous latent space compatible with Large Language Models. Unlike traditional raster-based approaches (BLIP-2, Florence-2), which render SVGs to pixel grids and process them with vision encoders, my method preserves the vector nature of the data and leverages the symbolic reasoning capabilities of LLMs.

The core methodology involves three stages: (1) SVG preprocessing and linearization, where the XML scene graph is traversed and path commands are extracted and normalized; (2) LoRA-based fine-tuning of the LLM, initially established as a text-only baseline to adapt the model to SVG syntax; and (3) geometric embedding via the SPE, which projects continuous coordinates through sinusoidal features to enable the full multimodal approach.

Experimentally, I evaluate the system on a dataset of 90,000 SVG-caption pairs from Icons8, with a stratified test set of 400 samples. I compare zero-shot raster baselines (BLIP-2, Florence-2), the \textbf{Baseline Training Process (LoRA-LLM)} (text-only fine-tuning on SVG code), and the \textbf{Multimodal Training Process (SPE-LoRA)} (my proposed vector-native approach). The results demonstrate that the \textbf{SPE+Qwen2-7B} configuration achieves the highest composite score (\textbf{8.89}, calculated as $\text{CLIPScore}/10 + \text{BLEU} + \text{METEOR} + \text{ROUGE}$), outperforming both zero-shot baselines and the text-only baseline in linguistic quality metrics (BLEU-1: \textbf{0.420}, METEOR: \textbf{0.380}, ROUGE-L: \textbf{0.450}), while maintaining competitive visual-semantic alignment (CLIPScore: \textbf{29.30}).

Qualitative analysis reveals that the model excels at describing geometric primitives and capturing spatial relationships in well-defined icon categories (User Interface, Arrows, Weather), but struggles with abstract shapes, color attribution (when style features are not encoded), and text recognition in path-based glyphs. These findings highlight the importance of feature normalization in the SPE pipeline and suggest directions for future work, including hierarchical encoding of SVG group structures and integration of style attributes.

\section{Thesis Organization}
\label{sec:thesis-organization}



The remainder of this manuscript is organized as follows. \textbf{Chapter 2} provides a comprehensive survey of the theoretical foundations and related work, moving from deep learning fundamentals to vision-language models and the specific technologies underpinning this thesis. \textbf{Chapter 3} presents the proposed methodology in detail, describing the design rationale for the SVG Path Embedder, the LoRA-based fine-tuning strategy, and the geometric intuition behind the architecture. \textbf{Chapter 4} documents the implementation, covering the software architecture, data loading pipeline, and training workflow. \textbf{Chapter 5} presents the experimental evaluation, reporting quantitative results for all model configurations and providing a detailed qualitative analysis of case studies and failure modes. Finally, \textbf{Chapter 6} synthesizes the findings, discusses limitations and ethical considerations, and outlines promising directions for future research. In the following chapters, I develop each of these themes in depth, building a rigorous and comprehensive account of the SVG captioning problem, my proposed solution, and the empirical evidence supporting its effectiveness.

\chapter{State of the Art}
\label{chap:stato-arte}

This chapter surveys the theoretical foundations and prior research underpinning the SVG captioning system. I move from general deep learning principles to specific technologies employed in my architecture, critically assessing their relevance to vector graphics understanding. The chapter begins with Transformer fundamentals (Section~\ref{sec:dl-foundations}), traces the evolution of Large Language Models (Section~\ref{sec:llms-evolution}), examines vision-language models and their limitations for vector data (Section~\ref{sec:vision-language}), explores positional encodings as the foundation for my SPE design (Section~\ref{sec:positional-encodings}), compares parameter-efficient fine-tuning methods (Section~\ref{sec:peft-methods}), and reviews prior work on SVG generation (Section~\ref{sec:vector-graphics-generation}).

\section{Foundations: Transformers and Self-Attention}
\label{sec:dl-foundations}

\subsection{The Transformer Architecture}

The Transformer \cite{vaswani2017attention} represented a paradigm shift in deep learning, replacing recurrent architectures with \textbf{self-attention} to enable parallel processing and efficient modeling of long-range dependencies. Given an input sequence of $N$ tokens represented as $d$-dimensional vectors, the mechanism constructs three matrices via learned linear projections: $Q = XW_Q$ (queries), $K = XW_K$ (keys), and $V = XW_V$ (values). Attention computes weighted sums:
\[
\text{Attention}(Q, K, V) = \text{softmax}\left(\frac{QK^T}{\sqrt{d_k}}\right)V
\]
Intuitively, the product $QK^T$ measures pairwise token similarities, the softmax normalizes these to probabilities, and the multiplication by $V$ aggregates relevant information. This mechanism is permutation-invariant, necessitating the use of positional encodings (discussed in Section~\ref{sec:positional-encodings}) to retain sequence order.

\begin{figure}[ht]
    \centering
    \includegraphics[width=0.8\textwidth]{figures/transformer_attention_is_all_you_need.png}
    \caption{The Transformer model architecture. Left: Encoder. Right: Decoder. (Source: Vaswani \textit{et al.}, 2017)}
    \label{fig:transformer-architecture}
\end{figure}


To capture different aspects of the input simultaneously, the architecture employs \textbf{Multi-head attention}, which projects the input into $h$ subspaces:
\[
\text{MultiHead}(Q, K, V) = \text{Concat}(\text{head}_1, \ldots, \text{head}_h)W_O
\]
Different heads can thus focus on distinct patterns, such as short-range syntactic dependencies or long-range semantic relationships (\textit{e.g.}, matching XML tags). Each Transformer layer combines this multi-head attention with a position-wise feed-forward network (FFN), defined as $\text{FFN}(x) = \text{GELU}(xW_1 + b_1)W_2 + b_2$, and wraps these sub-layers with residual connections and layer normalization: $\text{Output} = \text{LayerNorm}(x + \text{SubLayer}(x))$.

Modern LLMs (Qwen2, Gemma) use \textbf{RMSNorm} \cite{zhang2019root}, which omits mean centering for efficiency:
\[
\text{RMSNorm}(x) = \frac{x}{\text{RMS}(x)} \cdot \gamma, \quad \text{RMS}(x) = \sqrt{\frac{1}{d}\sum_{i=1}^d x_i^2}
\]

\subsection{SVG and the Tokenization Challenge}

SVG poses a unique challenge: while syntactically XML, its semantics are encoded in \textit{numeric coordinates}. A text tokenizer splits \texttt{M 12.5 34.7 L 56.2 89.3} into \texttt{[M, 12, ., 5, 34, ., 7, ...]}, destroying coordinate unity and losing numerical continuity. A minor change ($12.5 \to 12.6$) drastically alters the rendered shape, yet the tokenizer treats these as independent tokens. This fundamental mismatch motivates the design of the SVG Path Embedder--a specialized encoder preserving the continuous geometric nature of the data, with its theoretical foundation in Fourier-based positional encodings extended to 2D spatial data (Section~\ref{sec:positional-encodings}).

\section{Large Language Models: Evolution to Decoder-Only Architectures}
\label{sec:llms-evolution}

\subsection{From BERT to GPT: Encoder vs. Decoder}

The fundamental divergence in Transformer architectures stems from their training objectives and masking strategies, which dictate their suitability for different classes of tasks.

\textbf{BERT (Bidirectional Encoder Representations from Transformers)} \cite{devlin2018bert} employs a \textbf{Masked Language Modeling (MLM)} objective. In this paradigm, approximately 15\% of the input tokens are randomly replaced with a special \texttt{[MASK]} token (80\%), a random token (10\%), or left unchanged (10\%). The model trains to predict the original identity of these masked tokens based on the bidirectional context--that is, attending to both preceding and succeeding tokens simultaneously. Mathematically, the attention mechanism in BERT is unconstrained: the attention mask $M$ is a matrix of zeros, allowing
\[
A_{ij} = \text{softmax}\left(\frac{Q_i K_j^T}{\sqrt{d_k}}\right)
\]
to be non-zero for all pairs $(i, j)$ in the full sequence length $T$. This architecture builds deep, context-aware representations ideal for \textit{Natural Language Understanding (NLU)} tasks such as classification, named entity recognition, and question answering, where the entire input is available at once. However, this reliance on bidirectional context makes it inherently unsuitable for autoregressive generation, as the generation process is fundamentally inhibited by the look-ahead bias inherent in bidirectional models.

Conversely, \textbf{GPT (Generative Pre-trained Transformer)} \cite{radford2018improving, brown2020language} is designed for \textbf{Causal Language Modeling (CLM)}. It utilizes a decoder-only architecture where self-attention is constrained by a causal mask $M$, defined as:
\[
M_{ij} = \begin{cases} -\infty & \text{if } j > i \\ 0 & \text{otherwise} \end{cases}
\]
This ensures that the prediction of token $x_t$ depends exclusively on the history $x_{<t}$. The joint probability of a sequence is factorized as:
\[
P(x_1, \ldots, x_T) = \prod_{t=1}^T P(x_t \mid x_{<t})
\]
This factorization aligns perfectly with the generation process. While early GPT models lagged behind BERT in NLU benchmarks, scaling laws \cite{kaplan2020scaling} revealed that sufficiently large decoder-only models (starting with GPT-3) develop emergent capabilities--such as \textbf{In-Context Learning} and few-shot reasoning--that rival or surpass specialized encoder models. Unlike BERT, which typically requires fine-tuning a specific classification head for each downstream task, GPT models can perform tasks essentially "zero-shot" by continuing a prompt. For SVG captioning, this generative capability is paramount, as we must produce a sequential textual description from a visual prefix.

\subsection{Modern Decoder-Only LLMs}

While the original GPT architecture laid the foundation, modern open-weights LLMs have introduced critical architectural refinements to improve stability, efficiency, and context handling. Three state-of-the-art backbones are utilized in this thesis, each embodying specific design philosophies:

\begin{itemize}
    \item \textbf{Llama (Meta)} \cite{touvron2023llama}: Llama 3 maintains the robust design choices that characterize the Llama family. A key feature is the use of \textbf{Rotary Positional Embeddings (RoPE)} \cite{su2024roformer} instead of absolute positional encodings. RoPE encodes position by rotating the query and key vectors in the complex plane:
    \[
    f_{q,k}(x_m, m) = R^d_{\Theta, m} x_m
    \]
    This formulation allows the attention score to depend only on the relative distance $m-n$, improving generalization to sequence lengths exceeding the training window--essential for processing verbose SVG code. Additionally, it adopts \textbf{SwiGLU} activations \cite{shazeer2020glu} in the feed-forward networks, replacing ReLU with a gated mechanism $\text{Swish}(xW) \otimes (xV)$, which offers smoother gradients and faster convergence.
    
    \item \textbf{Qwen2 (Alibaba)} \cite{qwen2023}: The Qwen2-7B model is optimized for efficient long-context inference. It employs \textbf{Grouped-Query Attention (GQA)} \cite{ainslie2023gqa}, which serves as an effective trade-off between the speed of Multi-Query Attention (MQA) and the quality of Multi-Head Attention (MHA). In GQA, multiple query heads share a single key-value head (\textit{e.g.}, 8 KV heads for 64 Q heads). This drastically reduces the size of the KV cache and the memory bandwidth required during decoding, enabling the processing of context windows up to 128k tokens with minimal performance degradation.
    
    \item \textbf{Gemma (Google)} \cite{gemma2024}: Built on Gemini technology, Gemma focuses on numerical stability at scale. It places RMSNorm \textit{after} the embedding layer and before each transformer block (Pre-Norm with offset), preventing gradient explosions in deep networks. It also utilizes \textbf{GeGLU} activations (a Gaussian Error Linear Unit variant of GLU), which provides a more complex non-linear mapping than GeLU, enhancing the model's capacity to learn subtle patterns in data.
\end{itemize}

\textbf{Scaling Laws and Emergent Capabilities}.
Optimizing these architectures requires adhering to empirical scaling laws. While Kaplan \textit{et al.} \cite{kaplan2020scaling} initially suggested a power-law relationship favoring model size, the "Chinchilla" scaling laws \cite{hoffmann2022training} revised this, demonstrating that optimal performance for a given compute budget is achieved by scaling model parameters ($N$) and training tokens ($D$) in equal proportion ($N \approx D$). This insight drives the recent trend of highly capable "smaller" models (7B-9B parameters) trained on massive datasets (trillions of tokens). Importantly, as these models scale in accordance with these laws, they acquire \textit{emergent capabilities}--qualitative leaps in performance not predictable from smaller scales. For our specific multimodal task, the relevant emergence is the ability to map abstract vector structures (encoded by the SPE) to semantic linguistic concepts without requiring a dedicated cross-attention mechanism, effectively treating the SVG sequence as a native foreign language.



\section{Vision-Language Models and Their Limitations}
\label{sec:vision-language}
\subsection{Vision Transformers (ViT)}

The Vision Transformer (ViT) \cite{dosovitskiy2020image} marked a paradigm shift in computer vision, demonstrating that the pure Transformer architecture--without any convolutional inductive biases--could achieve state-of-the-art performance on image recognition tasks. ViT processes an image by dividing it into a grid of fixed-size patches (\textit{e.g.}, $16 \times 16$ pixels). Each patch is flattened into a 1D vector and linearly projected to the model's embedding dimension $D$. To retain spatial information, learnable position embeddings are added to these patch embeddings. A special learnable \texttt{[CLS]} token is prepended to the sequence, whose state at the output of the Transformer encoder serves as the image representation.

\begin{figure}[h]
    \centering
    \includegraphics[width=0.9\textwidth]{figures/architecture.png}
    \caption{The Vision Transformer (ViT) Architecture. The image is split into patches, linearly projected, and processed by a standard Transformer encoder. (Source: Dosovitskiy \textit{et al.}, 2020)}
    \label{fig:vit-architecture}
\end{figure}

Mathematically, if an input image $x \in \mathbb{R}^{H \times W \times C}$ is split into $N$ patches $x_p \in \mathbb{R}^{N \times (P^2 \cdot C)}$, the input sequence $z_0$ is:
\[
z_0 = [x_{class}; x_p^1 E; x_p^2 E; \ldots; x_p^N E] + E_{pos}
\]
where $E \in \mathbb{R}^{(P^2 \cdot C) \times D}$ is the projection matrix and $E_{pos} \in \mathbb{R}^{(N+1) \times D}$ are the position embeddings.

\subsection{CLIP: Contrastive Language-Image Pre-training}

CLIP \cite{radford2021learning} builds on ViT to learn visual representations from natural language supervision. Instead of training on fixed classes (like ImageNet), CLIP is trained on a massive dataset of 400 million image-text pairs collected from the internet. The core mechanism is \textbf{contrastive learning}. CLIP trains two encoders jointly: a vision encoder (ViT or ResNet) and a text encoder (Transformer). The goal is to maximize the cosine similarity between the embeddings of matched image-text pairs in a batch, while minimizing the similarity for unmatched pairs.

\begin{figure}[h]
    \centering
    \includegraphics[width=1.0\textwidth]{figures/clip_architecture_final.png}
    \caption{Architecture of CLIP (Contrastive Language-Image Pre-training). The model learns to align raster images and text in a shared embedding space, enabling zero-shot transfer to downstream tasks.}
    \label{fig:clip-architecture}
\end{figure}

This objective aligns the visual and textual latent spaces, enabling \textbf{zero-shot classification}: to classify an image, the model compares its embedding with the embeddings of text prompts like "a photo of a dog", "a photo of a cat", etc., and selects the highest match. However, CLIP operates fundamentally on raster inputs. Applying it to vector graphics requires rasterization, which discards the structural information inherent in the SVG code.

\subsection{Generative Captioning: BLIP-2 and LLaVA}

BLIP-2 \cite{li2023blip2} freezes a vision encoder (CLIP) and an LLM, training a lightweight Q-Former to bridge them. LLaVA \cite{liu2023llava} fine-tunes an LLM (Vicuna/Llama) with visual instruction data, demonstrating complex visual reasoning. While powerful, these models require raster inputs. My experiments (Chapter~\ref{chap:valutazione}) use BLIP-2 as a zero-shot baseline (rasterizing SVGs), revealing generic captions (``an icon'') lacking geometric precision.

\subsection{The Vector Graphics Gap}

Vision-language models designed for natural images have fundamental limitations when applied to SVG data. Specifically, they (i) require rasterization, which discards the structural information inherent in the vector format; (ii) struggle with resolution-dependent details, such as thin lines that may vanish at lower resolutions; and (iii) cannot leverage the explicit geometric relationships present in the XML markup. This gap motivates the development of the SPE--a vector-native encoder capable of processing SVG as structured, continuous geometric data.

These limitations are not merely theoretical but practical. Rasterization introduces a fixed resolution ceiling, causing "staircase" artifacts on diagonal lines and blurring fine details like small text or intricate icons. Furthermore, the raster representation is opaque; the model sees a grid of pixels but loses access to the semantic hierarchy (groups, layers, object types) explicitly defined in the SVG DOM. This forces the model to relearn geometric primitives from scratch, a task that is computationally inefficient and prone to error, especially for complex technical diagrams where precision is paramount. By treating SVG as a first-class citizen--ingesting the vector code directly--we can bypass these bottlenecks and provide the LLM with a lossless, semantically rich representation of the visual content.

\section{Positional and Coordinate Encodings}
\label{sec:positional-encodings}

Designing the SPE requires representing continuous 2D coordinates $(x, y)$ as dense vectors. This section traces the theory from NLP positional encodings to neural rendering.

In natural language processing, positional encodings are used to inject order information into the permutation-invariant self-attention mechanism. However, for 2D vector graphics, the challenge is two-fold: we must encode not just the sequence order of the path commands, but the absolute spatial position of the control points on the 2D canvas. A naive approach of normalizing coordinates to $[0, 1]$ and feeding them as scalar values to an MLP fails due to the "spectral bias" of neural networks, which struggle to learn high-frequency functions (\textit{i.e.}, sharp edges and fine details) from low-dimensional inputs. To overcome this, we draw inspiration from the field of Neural Radiance Fields (NeRF), where high-frequency Fourier features are used to map low-dimensional coordinates to a higher-dimensional space, enabling the network to capture intricate geometric structures.

\subsection{Sinusoidal Encodings in Transformers}

The original Transformer \cite{vaswani2017attention} uses fixed sinusoidal encodings to inject order information into the permutation-invariant self-attention mechanism:
\begin{equation}
\label{eq:sinusoidal-encoding}
\text{PE}_{(pos, 2i)} = \sin\left(\frac{pos}{10000^{2i/d}}\right), \quad \text{PE}_{(pos, 2i+1)} = \cos\left(\frac{pos}{10000^{2i/d}}\right)
\end{equation}

From a signal processing view, this is a \textit{Fourier feature mapping}. Each sine-cosine pair corresponds to a frequency band: low frequencies capture global structure; high frequencies capture fine details. This allows the model to distinguish relative positions in the sequence. In the context of SVG processing via Large Language Models, coordinates are typically serialized as text tokens (\textit{e.g.}, "1", "0", ".", "5"). Consequently, they are subject to these standard 1D positional encodings based on their token index in the sequence, treating geometric data purely as a linguistic sequence.

\begin{figure}[h]
    \centering
    \includegraphics[width=0.6\textwidth]{figures/positional_encoding.png}
    \caption{Visualization of 2D Sinusoidal Positional Encodings. High-frequency components capture fine details, while low-frequency components encode global position, enabling the model to "see" the continuous geometry.}
    \label{fig:2d-pe}
\end{figure}

\section{Parameter-Efficient Fine-Tuning (PEFT)}
\label{sec:peft-methods}

Fine-tuning billion-parameter LLMs is expensive (memory, compute) and risks catastrophic forgetting. Full fine-tuning of Qwen2-7B in FP16 requires ~56 GB GPU memory (14 GB weights + 14 GB gradients + 28 GB optimizer states). Even with sufficient memory, overfitting to the new task risks losing general capabilities. PEFT methods update only a small parameter subset to mitigate these issues. The core idea behind Parameter-Efficient Fine-Tuning is to freeze the vast majority of the pre-trained model's parameters and inject a small number of trainable weights. This hypothesis relies on the observation that the "intrinsic dimension" of the adaptation task is low--meaning that the changes required to adapt a general-purpose language model to a specific downstream task (like SVG captioning) can be captured in a low-rank subspace. This approach not only drastically reduces the memory footprint during training (since gradients are only computed for the small subset of parameters) but also facilitates the storage and deployment of multiple task-specific adapters on top of a single frozen backbone.

\subsection{PEFT Methods Overview}

Several PEFT strategies exist, each with distinct trade-offs. \textbf{Adapter Layers} \cite{houlsby2019parameter} insert bottleneck modules between Transformer layers; while effective, they add inference latency due to the sequential computation. \textbf{Prefix Tuning} \cite{li2021prefix} optimizes soft prompts prepended to the input, which is extremely parameter-efficient but consumes valuable context window space--a critical drawback for long SVG sequences.

In contrast, \textbf{Low-Rank Adaptation (LoRA)} \cite{hu2021lora} represents weight updates as low-rank products $\Delta W = B A$, where $B \in \mathbb{R}^{d \times r}$ and $A \in \mathbb{R}^{r \times k}$ with $r \ll \min(d, k)$. During training, the pre-trained weights $W_0$ are frozen, and only $A$ and $B$ are updated. At inference, the weights are merged ($W_{\text{merged}} = W_0 + \frac{\alpha}{r} B A$), yielding zero computational overhead.

For SVG captioning, LoRA is the optimal choice: it introduces no inference latency, preserves the full context window for long SVG-derived sequences, and offers proven training stability. This allows us to adapt decoder-only LLMs to the geometric modality on standard hardware while largely retaining their pre-trained linguistic capabilities.



\section{Vector Graphics Generation and Understanding}
\label{sec:vector-graphics-generation}

Research on vector graphics has accelerated dramatically in recent years, driven by the desire to extend the success of deep generative models from the raster domain to the vector domain. Unlike pixels, vector graphics are non-grid-structured, variable-length, and symbolic, requiring specialized architectures. I categorize the state-of-the-art into four primary families: Autoregressive/LLM-based approaches, Diffusion-based methods, Representation Learning, and Parameter Optimization.

This line of work treats SVG generation as a sequence modeling problem, leveraging the sequential nature of the SVG file format (XML text or path commands). Wu \textit{et al.} introduced \textbf{IconShop} \cite{wu2023iconshop}, which represents a significant milestone in autoregressive SVG generation. It utilizes a Transformer-based decoder to generate SVG primitives sequentially, conditioned on text prompts. The key innovation lies in its "autoregressive with instruction tuning" strategy, where the model learns not just to complete a sequence, but to follow high-level design constraints. By tokenizing the SVG commands (MoveTo, LineTo, CubicBezier) into a discrete vocabulary, IconShop frames vector generation as a language modeling task.

\begin{figure}[htbp]
    \centering
    \includegraphics[width=0.9\textwidth]{figures/screenshot_3.png}
    \caption{Overview of Autoregressive SVG Generation (IconShop). The model predicts the next drawing command based on the history of previous commands and a text condition. (Source: Wu \textit{et al.}, 2023)}
    \label{fig:iconshop-arch}
\end{figure}
\FloatBarrier

Building on this, \textbf{StarVector} \cite{rodriguez2023starvector} introduces a multimodal code-visual architecture. It recognizes that SVGs are both code (XML) and visual content. StarVector employs a dual-encoder setup: a code encoder processes the raw SVG tokens, while a visual encoder (like ViT) processes the rendered image. These representations are fused to guide the generation process, allowing the model to understand complex spatial relationships that are implicit in the code but explicit in the rendering.

\begin{figure}[htbp]
    \centering
    \includegraphics[width=1.0\textwidth]{figures/screenshot_2.png}
    \caption{StarVector Architecture. A dual-encoder system that processes both the raw SVG code and the rendered image to guide generation. (Source: Rodriguez \textit{et al.}, 2023)}
    \label{fig:starvector-arch}
\end{figure}
\FloatBarrier

Most recently, Zini \textit{et al.} \cite{zini2025vhector} introduced \textbf{vHector}, a system designed to improve the scalability and quality of vector generation. vHector linearizes the SVG DOM tree into a sequence of tokens but annotates them with rich geometric metadata, such as absolute/relative coordinate flags and styling cues. This "augmented tokenization" helps the LLM maintain spatial coherence over long sequences. Furthermore, the accompanying \textbf{HeisenVec} module acts as a denoising autoencoder for vector primitives, refining the output of the LLM to ensure valid and aesthetically pleasing shapes.
\FloatBarrier

Inspired by the success of Stable Diffusion and DALL-E 2, these methods adapt probabilistic diffusion processes to the vector domain. \textbf{VectorFusion} \cite{jain2023vectorfusion} was one of the first to apply \textit{Score Distillation Sampling (SDS)} to vector graphics. Instead of training a diffusion model on SVGs directly (which is hard due to the lack of large-scale vector datasets), VectorFusion optimizes a set of vector paths to match a text prompt using a pre-trained frozen image diffusion model (like Stable Diffusion) as a critic. The gradients from the image diffusion model are backpropagated through a differentiable rasterizer (DiffVG) to update the SVG parameters.

\begin{figure}[htbp]
    \centering
    \includegraphics[width=1.0\textwidth]{figures/screenshot_1.png}
    \caption{VectorFusion Pipeline. A set of random paths is optimized via Score Distillation Sampling (SDS) from a frozen text-to-image diffusion model. The differentiable rasterizer allows gradients to flow back to the vector parameters. (Source: Jain \textit{et al.}, 2023)}
    \label{fig:vectorfusion-arch}
\end{figure}
\FloatBarrier

\textbf{SVGDreamer} \cite{xing2024svgdreamer} improves upon VectorFusion by addressing the "initialization problem." VectorFusion often gets stuck in local minima if initialized randomly. SVGDreamer introduces a two-stage process: first, it generates a semantic attention map to guide the placement of paths; second, it optimizes the paths using a specialized "Vectorized Score Distillation" loss that encourages cleaner geometry and better separation of foreground/background elements.
\FloatBarrier

This category focuses on learning a latent space for vector graphics that captures their structural and hierarchical properties, enabling tasks like interpolation, classification, and retrieval. \textbf{DeepSVG} \cite{carlier2020deepsvg} is a foundational work in this area. It proposes a hierarchical Transformer architecture that decomposes an SVG into a hierarchy of paths and commands. The encoder compresses the SVG into a global latent vector, while the decoder reconstructs it command-by-command. This hierarchical approach allows DeepSVG to model complex multi-path icons effectively and perform latent space arithmetic (\textit{e.g.}, morphing one icon into another).

\textbf{SuperSVG} \cite{hu2024supersvg} introduces a "superpixel" concept to vector learning. It decomposes an image into superpixels (clusters of pixels with similar color/texture) and models the relationships between them using a Graph Neural Network (GNN). This graph representation is then converted into vector paths. This approach is particularly effective for vectorizing complex natural images where standard path tracing fails.
\FloatBarrier

Unlike generative models that predict new geometry, optimization approaches start with a canvas of shapes and adjust their parameters (coordinates, colors) to minimize a loss function. \textbf{LIVE} (Layer-wise Image Vectorization) \cite{ma2022live} focuses on converting raster images to SVG. It initializes a set of Bézier curves and optimizes their control points to minimize the reconstruction error against the target image. The key enabler is \textbf{DiffVG} \cite{li2020differentiable}, a differentiable rasterizer that allows gradients to flow from pixel-space loss to vector-space parameters.

\textbf{ClipVG} \cite{song2023clipvg} combines differentiable rasterization with CLIP. It optimizes the parameters of SVG paths so that the rendered image maximizes the CLIP similarity with a target text prompt. This allows for "text-to-SVG" generation without a generative model, purely through optimization. However, this process is slow and often results in "messy" geometry (many overlapping paths) compared to autoregressive methods.



\input{chapter3_expanded}

% ============================================
% BIBLIOGRAFIA
% ============================================

\bibliographystyle{plain}
\bibliography{bibliography}

\end{document}
